\clearpage\fancyhead[L]{\small PALLAS. KOMET.}\hypertarget{native-332-lower}{}\label{native-332-lower}
\begin{center}\textbf{PALLAS. KOMET.}\par\vspace{4pt}\hrule\vspace{4pt}\fontsize{10}{12}\selectfont Göttingische gelehrte Anzeigen. Stück 130. Seite 1289\,..\,1293. 1811. August 17.\par\vspace{4pt}\hrule\vspace{4pt}\end{center}

\par\indent Im 8. Stück unserer Anzeigen dieses Jahres [S. 322 d. B.] wurde über die Arbeit von Professor \textsc{Gauss} zu den Störungen der Pallas seit ihrer Entdeckung berichtet. Dort steht auch die \SourceLeaf{333}erste Beobachtung vom 4. Januar an unserer Sternwarte, die die Ergebnisse dieser Arbeit mit einer kaum erwarteten Genauigkeit bestätigte. Der Planet war zu lichtschwach, um an unserer Sternwarte im Meridian beobachtet werden zu können. Die Beobachtungen wurden daher hier nicht fortgesetzt. Lediglich auswärtige Astronomen, denen leistungsfähigere Instrumente zur Verfügung standen, wurden rechtzeitig benachrichtigt. So erhielt Professor \textsc{Gauss} für die Zeit der Opposition, die im Februar eintrat, sehr wertvolle Beobachtungen von fünf Sternwarten, nämlich aus Paris, Mannheim, Seeberg, Berlin und Hamburg. Sie dienten ihm als Grundlage für seine weiteren Untersuchungen. Die Beobachtungen der Herren \textsc{von Lindenau}, \textsc{Barry} und \textsc{Harding} an den Sternwarten auf dem Seeberg und in Mannheim sind bereits in der Monatlichen Correspondenz abgedruckt. Die übrigen teilen wir hier mit.

\Needspace{14\baselineskip}\begin{center}\fontsize{10}{12}\selectfont\setlength{\tabcolsep}{4pt}\renewcommand{\arraystretch}{1.15}
Beobachtungen der Pallas in Paris, von Herrn \textsc{Bouvard}.

\begin{tabular}{c|c|c}
im Meridian & Rektaszension & Südliche Deklination \\
\hline
1811 Febr. 17 & $146^\circ\,58'\,33''15$ & $13^\circ\,29'\,42''4$ \\
18 & $146\;\;47\;\;55.95$ & $13\;\;\;5\;\;59.4$ \\
22 & $146\;\;\;6\;\;27.90$ & $11\;\;28\;\;19.6$ \\
27 & $145\;\;18\;\;13.65$ & $9\;\;19\;\;32.9$ \\
\end{tabular}
\end{center}

\Needspace{14\baselineskip}\begin{center}\fontsize{10}{12}\selectfont\setlength{\tabcolsep}{4pt}\renewcommand{\arraystretch}{1.15}
Beobachtungen in Berlin, von Professor \textsc{Bode}.

\begin{tabular}{c|c|c|c}
1811 & Mittlere Zeit & Rektaszension & Südliche Deklination \\
\hline
Febr. 18 & $11^{\mathrm u}\,55^{\mathrm m}\,22^{\mathrm s}5$ & $146^\circ\,48'\,2''$ & --- --- --- \\
19 & $11\;\;50\;\;34.0$ & $146\;\;37\;\;32$ & $12^\circ\,42'\,44''5$ \\
20 & $11\;\;45\;\;57.0$ & $146\;\;27\;\;\;4$ & $12\;\;19\;\;\;1$ \\
März 12 & $10\;\;16\;\;33.6$ & $143\;\;45\;\;29$ & $3\;\;32\;\;42$ \\
16 & $9\;\;59\;\;45.1$ & $143\;\;29\;\;18$ & $1\;\;48\;\;\;3$ \\
17 & $9\;\;55\;\;35.1$ & $143\;\;26\;\;13$ & $1\;\;22\;\;11$ \\
\end{tabular}
\end{center}

\Needspace{14\baselineskip}\begin{center}\fontsize{10}{12}\selectfont\setlength{\tabcolsep}{4pt}\renewcommand{\arraystretch}{1.15}
Beobachtungen in Hamburg, von Professor \textsc{Schumacher}.

\begin{tabular}{c|c|c}
im Meridian & Rektaszension & Südliche Deklination \\
\hline
1811 Febr. 19 & $146^\circ\,37'\,15''0$ & $12^\circ\,42'\,45''3$ \\
20 & $146\;\;26\;\;52.5$ & $12\;\;18\;\;32.7$ \\
21 & $146\;\;16\;\;37.5$ & $11\;\;53\;\;47.7$ \\
22 & $146\;\;\;6\;\;36.0$ & $11\;\;28\;\;53.5$ \\
März 9 & $144\;\;\;1\;\;45.0$ & $4\;\;52\;\;42.5$ \\
15 & $143\;\;32\;\;39.0$ & $2\;\;13\;\;59.0$ \\
16 & $143\;\;29\;\;24.0$ & $1\;\;47\;\;57.0$ \\
18 & $143\;\;23\;\;34.5$ & $0\;\;56\;\;47.8$ \\
\end{tabular}
\end{center}


\Needspace{8\baselineskip}\noindent Daraus leitete Professor \textsc{Gauss} das folgende Ergebnis für die Opposition ab.
\begin{center}1811 Febr. 21. $19^{\mathrm u}\,25^{\mathrm m}\,31^{\mathrm s}$ mittlere Zeit in Göttingen\end{center}

\begin{center}\fontsize{10}{12}\selectfont\setlength{\tabcolsep}{4pt}\renewcommand{\arraystretch}{1.15}\begin{tabular}{llr}
Wahre heliozentrische Länge & .\;.\;.\;.\;.\;.\;.\;. & $152^\circ\,48'\,15''8$\\
Wahre geozentrische Breite, südlich & .\;.\;.\;.\;. & $23\;48\;19.2$\\
\end{tabular}\end{center}

\Needspace{17\baselineskip}\par\indent Wenn man die von Professor \textsc{Gauss} berechneten Störungen durch Jupiter ordnungsgemäß berücksichtigt, ist die Übereinstimmung dieser Opposition mit den sechs vorhergehenden in der Tat bewundernswert. Nachdem an den schon vor der Wiederauffindung bestimmten Elementen nur einige äußerst geringfügige Korrekturen vorgenommen worden waren, ergaben sich die folgenden Unterschiede zwischen Beobachtung und Rechnung.
\SourceLeaf{334}
\begin{center}\fontsize{10}{12}\selectfont\setlength{\tabcolsep}{4pt}\renewcommand{\arraystretch}{1.15}
Unterschied der

\begin{tabular}{c|c|c}
 & mittleren Länge & heliozentrischen Breite \\
\hline
1803 & $+8''7$ & $-1''1$ \\
1804 & $-7.0$ & $-10.2$ \\
1805 & $+1.4$ & $-3.5$ \\
1807 & $-1.2$ & $-2.4$ \\
1808 & $-3.1$ & $-21.4$ \\
1809 & $+4.5$ & $-2.1$ \\
1811 & $-2.2$ & $-12.1$ \\
\end{tabular}
\end{center}


\Needspace{21\baselineskip}\par\indent Anschließend wurde die Berechnung der Störungen um ein weiteres Jahr fortgeführt. Professor \textsc{Gauss} wurde bei dieser mühsamen Arbeit von Herrn \textsc{Nicolai} unterstützt, der sich hier mit großem Eifer den mathematischen Wissenschaften widmet. Unter der Aufsicht von Professor \textsc{Gauss} führte er den größeren Teil der numerischen Rechnungen mit ebenso viel Fleiß wie Geschick aus. Nach den Ergebnissen dieser Rechnung wird sich die Bewegung der Pallas um die Zeit der Opposition im nächsten Jahr durch die folgenden elliptischen Elemente darstellen lassen, in denen die Störungen bereits enthalten sind.

\begin{center}\fontsize{10}{12}\selectfont\setlength{\tabcolsep}{4pt}\renewcommand{\arraystretch}{1.15}\begin{tabular}{llr}
\shortstack[l]{Epoche der mittleren Länge 1812 Juni 10 Mittag\\im Meridian von Göttingen} & . & $239^\circ\,4'\,46''1$\\
Länge des Perihels für dieselbe Zeit &  & $121\;0\;48.5$\\
Länge des aufsteigenden Knotens & . & $172\;32\;44.5$\\
Neigung der Bahn & .\;.\;.\;.\;.\;.\;. & $34^\circ\,34'\,54''7$\\
Tägliche mittlere tropische Bewegung &  & $768''5746$\\
Exzentrizität & .\;.\;.\;.\;.\;. & $=\sin13^\circ\,59'\,1''8$\\
Logarithmus der großen Halbachse & . & $0.4429321$\\
\end{tabular}\end{center}

\Needspace{20\baselineskip}\par\indent Aus eben diesen Elementen hat Herr \textsc{Nicolai} auch eine Ephemeride für die Erscheinung der Pallas im Jahr 1812 berechnet. Sie ist zu umfangreich, um hier Platz zu finden, und wird in der Monatl. Correspondenz und im Astronom. Jahrbuch abgedruckt werden. Wir geben nur noch die ebenfalls von Herrn \textsc{Nicolai} \emph{im Voraus} berechnete nächste Opposition (die achte) an. Sie kann zu gegebener Zeit dazu dienen, die Genauigkeit unserer Ergebnisse zu prüfen.
\begin{center}1812 Juni 10. $3^{\mathrm u}\,32^{\mathrm m}$ mittlere Zeit in Göttingen.\end{center}

\begin{center}\fontsize{10}{12}\selectfont\setlength{\tabcolsep}{4pt}\renewcommand{\arraystretch}{1.15}\begin{tabular}{llr}
Heliozentrische Länge der Pallas & .\;.\;.\;.\;.\;.\;.\;. & $259^\circ\,28'\,53''$\\
Heliozentrische Breite, nördlich & .\;.\;.\;.\;.\;.\;.\;. & $34\;32\;37$\\
Geozentrische Breite & .\;.\;.\;.\;.\;.\;.\;.\;.\;.\;.\;. & $48\;16\;36$\\
\end{tabular}\end{center}

\begin{center}\fontsize{10}{12}\selectfont\setlength{\tabcolsep}{4pt}\renewcommand{\arraystretch}{1.15}\begin{tabular}{lr}
\shortstack[l]{Logarithmus der Entfernung der Pallas\\von der Sonne} & $0.504259$\\
\shortstack[l]{Logarithmus der Entfernung der Pallas\\von der Erde} & $0.384915$\\
\end{tabular}\end{center}

\Needspace{20\baselineskip}\par\indent Bei dieser Gelegenheit fügen wir noch die parobolischen Elemente des Kometen hinzu, den \textsc{Flaugergues} im März dieses Jahres entdeckt hat. Professor \textsc{Gauss} hat sie soeben aus einigen Beobachtungen berechnet, die Herr von \textsc{Zach} in Marseille angestellt hat. Die Elemente sind vorläufig und können die Wiederauffindung dieses bemerkenswerten Kometen gegen Ende August erleichtern.
\begin{center}Durchgang durch das Perihel 1811 Sept. 10. $0^{\mathrm u}\,51^{\mathrm m}$\end{center}

\begin{center}\fontsize{10}{12}\selectfont\setlength{\tabcolsep}{4pt}\renewcommand{\arraystretch}{1.15}\begin{tabular}{llr}
Länge des Perihels & .\;.\;.\;.\;.\;.\;.\;.\;.\;.\;.\;. & $73^\circ\,14'\,35''$\\
Logarithmus des kleinsten Abstands & .\;.\;.\;.\;.\;.\;. & $9.99153$\\
Länge des aufsteigenden Knotens & .\;.\;.\;.\;.\;.\;.\;. & $141^\circ\,4'\,59''$\\
Neigung der Bahn & .\;.\;.\;.\;.\;.\;.\;.\;.\;.\;.\;.\;. & $73\;48\;2$\\
Bewegung \emph{rückläufig} &  & \\
\end{tabular}\end{center}

\SourceLeaf{335}\Needspace{12\baselineskip}\par\indent Aus diesen Elementen wurden die folgenden Positionen berechnet.
\begin{center}\fontsize{10}{12}\selectfont\setlength{\tabcolsep}{4pt}\renewcommand{\arraystretch}{1.15}
\begin{tabular}{l|cc|c}
Mitternacht & Rektaszension & Nördliche Deklination & Lichtstärke \\
\hline
1811 Aug. 11 & $142^\circ\,43'$ & $29^\circ\,54'$ & $0.194$ \\
\phantom{1811 }\phantom{Aug. }26 & $152\;\;19$ & $36\;\;15$ & $0.294$ \\
Sept. 10 & $167\;\;34$ & $43\;\;14$ & $0.434$ \\
\phantom{Sept. }25 & $194\;\;22$ & $48\;\;35$ & $0.576$ \\
Okt. 10 & $232\;\;28$ & $44\;\;52$ & $0.595$ \\
\end{tabular}
\end{center}


\par\indent Nach demselben Maßstab betrug die Helligkeit (unter der Voraussetzung, dass der Komet sein Licht allein von der Sonne erhält) am 11. April $0.035$ und am 2. Juni $0.052$. Anfang Oktober wird er also 17-mal so hell sein wie am 11. April und mit ganz besonderem Glanz am nördlichen Himmel erscheinen. Danach wird er bei abnehmender Helligkeit noch durch die Sternbilder Herkules und Adler ziehen und zu Beginn des nächsten Jahres im Wassermann unsichtbar werden. Wenn sein Licht bis dahin allmählich schwächer geworden ist, aber die Beobachtung noch zulässt, kann er im April vielleicht im Wassermann noch einmal sichtbar werden und mit sehr leistungsstarken Teleskopen noch eine geraume Zeit verfolgt werden. Aber auch wenn das nicht gelingen sollte, wird er wegen seines außerordentlichen Glanzes im Herbst dieses Jahres und wegen seiner langen Sichtbarkeit, die gewiss an die zehn Monate betragen wird, in den Annalen der Kometenkunde eine der allerbemerkenswertesten Erscheinungen bleiben.
